\chapter{Finite-Temperature Effective Potential in Cosmology}
\label{ch:fT-eff-pot-cosmo}

We now extend the analysis of the previous chapter to finite temperature.
The imaginary-time formalism developed in Chapter~\ref{ch:finite-temperature}
provides the framework: at temperature $T = 1/\beta$, the partition
function is a path integral over fields periodic in Euclidean time with
period $\beta$, and the loop integrals are replaced by Matsubara sums.
Here we apply this machinery to derive the finite-temperature effective
potential in the cosmological context, and use it to analyze the
temperature dependence of the phase structure.

\section{The Imaginary-Time Formalism in Brief}

The canonical ensemble in quantum mechanics is described by the density
matrix $\rho = Z^{-1}e^{-\beta H}$, with partition function
$Z = \mathrm{tr}(e^{-\beta H})$ and ensemble averages
\begin{equation}
	\langle\mathcal{O}\rangle
	= \mathrm{tr}(\rho\,\mathcal{O})
	= \frac{1}{Z}\sum_n e^{-\beta E_n}
	\langle\psi_n|\mathcal{O}|\psi_n\rangle.
\end{equation}
In quantum field theory the trace becomes a path integral over field
configurations,
\begin{equation}
	Z = \mathrm{tr}(e^{-\beta H})
	= \int\mathcal{D}\phi\,\langle\phi|e^{-\beta H}|\phi\rangle.
\end{equation}
To cast this in path integral form, we use Feynman's formula for the
transition amplitude,
\begin{equation}
	\langle\phi_b|e^{-iH(t_b-t_a)}|\phi_a\rangle
	= \int_{\phi(t_a)=\phi_a}^{\phi(t_b)=\phi_b}
	\mathcal{D}\phi\;
	e^{i\int_{t_a}^{t_b}dt\int d^3x\,\mathcal{L}},
\end{equation}
and analytically continue to imaginary time $\tau = it$. Setting
$t_a=0$ and $t_b = -i\beta$, identifying $\phi_a = \phi_b$ and
integrating over field configurations, the partition function becomes
\begin{equation}
	Z = \int_{\phi(0,\mathbf{x})=\phi(\beta,\mathbf{x})}
	\mathcal{D}\phi\;
	e^{-\int_0^\beta d\tau\int d^3x\,\mathcal{L}_E},
\end{equation}
where the Euclidean Lagrangian is
\begin{equation}
	\mathcal{L}_E = (\partial_\tau\phi)^2
	+ (\partial_i\phi)^2 + V(\phi).
\end{equation}
The fields satisfy periodic boundary conditions in Euclidean time,
$\phi(0,\mathbf{x}) = \phi(\beta,\mathbf{x})$ for bosons and
antiperiodic for fermions. The generating functional at finite
temperature is
\begin{equation}
	Z^\beta[J] = \int\mathcal{D}\phi\;
	e^{-\int_0^\beta d\tau\int d^3x\,
			[\mathcal{L}_E - J\phi]},
\end{equation}
and the chain from $Z^\beta[J]$ to the finite-temperature effective
potential $V_{\text{eff}}^\beta(\bar\phi_c)$ follows exactly the same
Legendre transform logic as at zero temperature,
\begin{equation}
	Z^\beta[J]
	\xrightarrow{\;\text{Legendre}\;}
	\Gamma^\beta[\bar\phi]
	\xrightarrow{\;\bar\phi=\bar\phi_c\;}
	V_{\text{eff}}^\beta(\bar\phi_c).
\end{equation}
The key difference from zero temperature lies in the Feynman rules. The
propagator becomes $-i/(\omega_k^2 + \mathbf{p}^2 + m^2)$ in Euclidean
space, the vertex is unchanged at $-i\lambda$, and the loop integral is
replaced by a Matsubara sum,
\begin{equation}
	\int\frac{d^4p}{(2\pi)^4}
	\;\longrightarrow\;
	\frac{i}{\beta}\sum_{k=-\infty}^{\infty}
	\int\frac{d^3p}{(2\pi)^3},
\end{equation}
with bosonic Matsubara frequencies $\omega_k = 2\pi k\beta^{-1}$ and
fermionic frequencies $\omega_k = (2k+1)\pi\beta^{-1}$.

\section{The Thermal Functions \texorpdfstring{$J_B$}{JB} and
  \texorpdfstring{$J_F$}{JF}}

Following the same steps as in Chapter~\ref{ch:finite-temperature},
the one-loop finite-temperature contribution to the effective potential
for a scalar field with field-dependent mass $m^2(\phi)$ is
\begin{equation}
	V_{1,\beta}(\phi) = \int\frac{d^3p}{(2\pi)^3}
	\left[\frac{\omega}{2}
		+ \frac{1}{\beta}\ln(1-e^{-\beta\omega})\right],
	\qquad \omega = \sqrt{\mathbf{p}^2+m^2(\phi)}.
\end{equation}
The first term is the zero-point energy, which upon Wick rotation
reproduces the zero-temperature Coleman-Weinberg result. The second
term is the genuinely thermal contribution. After the substitution
$x = |\mathbf{p}|\beta$, it takes the form
\begin{equation}
	V_{1,\beta}(\phi) = \frac{T^4}{2\pi^2}
	J_B\!\left(\frac{m^2(\phi)}{T^2}\right),
\end{equation}
where the bosonic thermal function is defined as
\begin{equation}
	J_B(y^2) = \int_0^{\infty}dx\,x^2
	\ln\!\left(1 - e^{-\sqrt{x^2+y^2}}\right).
\end{equation}
For a $\mathrm{U}(1)$ gauge field with field-dependent mass
$m_A^2(\phi) = 2g^2|\phi|^2$, the three physical polarizations
contribute
\begin{equation}
	V_{1,\beta}^{(A_\mu)}(\phi)
	= \frac{3T^4}{2\pi^2}
	J_B\!\left(\frac{m_A^2(\phi)}{T^2}\right).
\end{equation}
For Dirac fermions with field-dependent mass $m_f^2 = y^2\phi^2$,
the antiperiodic boundary conditions give fermionic Matsubara frequencies
$\omega_k = (2k+1)\pi\beta^{-1}$, and the thermal contribution is
\begin{equation}
	V_{1,\beta}^{(\psi)}(\phi)
	= -\frac{2T^4}{\pi^2}
	J_F\!\left(\frac{m_f^2(\phi)}{T^2}\right),
\end{equation}
where the fermionic thermal function is
\begin{equation}
	J_F(y^2) = \int_0^{\infty}dx\,x^2
	\ln\!\left(1 + e^{-\sqrt{x^2+y^2}}\right),
\end{equation}
with a plus sign in the exponential reflecting the antiperiodic boundary
conditions. The high-temperature expansions of these functions, valid
for $y^2 = m^2(\phi)/T^2\ll 1$, are
\begin{align}
	J_B(y^2) & \cong -\frac{\pi^4}{45}
	+ \frac{\pi^2}{12}y^2
	- \frac{\pi}{6}|y|^3
	- \frac{1}{32}y^4\ln\frac{y^2}{a_B}
	+ \cdots,                           \\
	J_F(y^2) & \cong \frac{7\pi^4}{360}
	- \frac{\pi^2}{24}y^2
	- \frac{1}{32}y^4\ln\frac{y^2}{a_F}
	+ \cdots,
\end{align}
where $a_B = \pi^2 e^{3/2-2\gamma_E}$ and
$a_F = 16\pi^2 e^{3/2-2\gamma_E}$. The leading terms $\mp\pi^4/45$
and $7\pi^4/360$ are the free energy densities of massless bosons and
fermions respectively, the Stefan-Boltzmann contributions. The
$\pm\pi^2 y^2/12$ terms are the temperature-dependent mass corrections
already computed in Chapter~\ref{ch:finite-temperature}. The cubic
term $-\pi|y|^3/6$ in $J_B$ is new and is absent for fermions: it
arises from the bosonic Matsubara zero mode and, as we shall see, plays
a crucial role in generating a potential barrier and hence a first-order
phase transition.

\section{The Coleman-Weinberg Model at Finite Temperature}

We now bring together the zero-temperature and finite-temperature
contributions to obtain the full effective potential of the
Coleman-Weinberg model at temperature $T$. The total effective potential
is
\begin{equation}
	V_{\text{eff}}(\phi,T)
	= \underbrace{V_0(\phi) + V_{1,\text{CW}}(\phi)}_{
		\lambda_{\text{eff}}(\phi)\phi^4/4!}
	+ V_{1,\beta}(\phi),
	\label{eq:tot-eff-pot}
\end{equation}
where the zero-temperature part was computed in the previous chapter.
In the regime $\lambda^2\ll g^4$, the thermal contribution is dominated
by the gauge boson loop,
\begin{equation}
	V_{1,\beta}(\phi)
	\cong \frac{3T^4}{2\pi^2}
	J_B\!\left(\frac{2g^2\phi^2}{T^2}\right).
\end{equation}
In the high-temperature approximation $T^2\gg m_A^2(\phi) = 2g^2\phi^2$,
expanding $J_B$,
\begin{equation}
	V_{1,\beta}(\phi)
	\cong \frac{g^2}{4\pi^2}T^2\phi^2
	+ \mathcal{O}(\phi^3) + \text{const.},
\end{equation}
so that the full effective potential becomes
\begin{equation}
	V_{\text{eff}}(\phi,T)
	= \frac{\lambda_{\text{eff}}(\phi)}{4!}\phi^4
	+ \frac{g^2}{4\pi^2}T^2\phi^2 + \cdots
\end{equation}
The positive $T^2\phi^2$ term drives symmetry restoration at high
temperature: it adds a positive mass-squared to the field, pushing the
minimum back toward $\phi = 0$. At sufficiently high $T$ the symmetric
phase is restored. As the universe cools, the coefficient of $\phi^2$
decreases and eventually the broken-symmetry minimum at $\phi\neq 0$
becomes energetically favored.

% [FIGURE: Effective potential at various temperatures showing symmetry
%  restoration at high T and the development of a minimum at low T]

The cubic term $\sim T|\phi|^3$ from $J_B$ is particularly significant.
It generates a potential barrier between the symmetric and broken-symmetry
phases that persists down to temperatures below the naive critical
temperature, making the transition first order. The height and width of
this barrier determine the nucleation rate of bubbles of the broken phase
within the symmetric phase, which we analyze in detail in
Part~\ref{part:false-vacuum}.

The temperature at which the two minima become degenerate is the
critical temperature $T_c$, derived in Chapter~\ref{ch:finite-temperature}
for the purely quartic case. Below $T_c$ the broken-symmetry phase is
the global minimum, and the universe undergoes the phase transition.
The order and character of this transition, and its cosmological
consequences, depend sensitively on the relative magnitudes of the
gauge and scalar couplings, and on the full structure of the effective
potential including the cubic term from the thermal bosonic loop.
